\section{A closed-interval survival theorem}
\label{sec:survival}

The estimates above apply before choosing a trajectory inside the
box. This allows confinement and survival to be proved together.
Write
\[
E_\beta(t)=\int_0^t e^{4\beta s}\,\dd s
=\begin{cases}
(e^{4\beta t}-1)/(4\beta),&\beta\ne0,\\
t,&\beta=0.
\end{cases}
\]
For every real \(\beta\), \(E_\beta\) increases with \(t\).

\begin{theorem}[Conditional finite-pair survival]
\label{thm:survival}
Let \(h\) be real symmetric, analytic in space and sufficiently
smooth in time for the heat equation and implicit continuation on
a neighborhood of the box. Suppose \(h_\tau=h_{zz}\) and that
\eqref{eq:deflation-error} and \eqref{eq:model-floor} hold on
\(\mathcal K\). At \(\tau=0\) let \(h\) have two specified distinct
simple real roots satisfying
\[
|c(0)-X|\le R_0<R,\qquad
0<Q_-\le q(0)\le Q_+<4D^2.
\]
Let \(\beta\) be an upper bound for \(B\) at each actual pair in the
box, for instance \(K\) or \(b_P+\kappa\) from Proposition~\ref{prop:force}.
If
\begin{equation}
\label{eq:survival-margins}
R_0+2A\Delta<R,\qquad
Q_+e^{4K\Delta}<4D^2,\qquad
\eta=Q_- -8E_\beta(\Delta)>0 ,
\end{equation}
then the pair is real, distinct and simple throughout
\([0,\Delta]\) and stays inside the spatial bounds of the box.
Its gap satisfies \(q(t)\ge e^{-4\beta t}(Q_- -8E_\beta(t))\).
\end{theorem}

\begin{proof}
Follow the initial roots by the implicit function theorem while
they are distinct and simple with \(|c-X|<R\), \(0<d<D\).
Real symmetry and uniqueness give real branches. The uniform
floor holds for the nodes on this initial interval of continuation,
so there is no third zero on the segment and \eqref{eq:gap-law}
applies. It gives \(|c'|\le2A\) and \(q'\le4Kq\), whence
\[
|c(t)-X|\le R_0+2At,\qquad q(t)\le Q_+e^{4Kt}.
\]
The lower inequality \(q'\ge -8-4\beta q\) gives
\(q(t)\ge e^{-4\beta t}(Q_- -8E_\beta(t))\) by an integrating
factor. In particular \(q(t)\ge e^{-4\max(\beta,0)\Delta}\eta>0\)
on the continuation interval. At either endpoint the slope obeys
\[
|h_z|\ge\sqrt q\,(m-e_0).
\]

If \(S\le\Delta\) is the first candidate stopping time, the bounds
on \(c'\) and \(q'\) give limits of \(c,q\) as \(t\uparrow S\).
Their associated root positions \(c\pm\sqrt q/2\) converge to
distinct zeros by continuity. Both slopes stay nonzero in the
limit and the spatial inequalities hold strictly. The implicit
function theorem extends the branches if \(S<\Delta\) and
establishes simplicity at \(S=\Delta\).
\end{proof}

An additional zero cannot merge with either endpoint during this
argument, since a merger would destroy its simple zero. This finite
real-pair conclusion requires no separate complex zero count.
The full virtual-node floor is a substantial input, and can fail
even when a real pair is regular. A failed sufficient inequality
leaves the fate of the actual pair undecided.

\subsection{Three elementary endpoint checks}

For \(h=(z-X)^2+2(\tau-T)\), exact modeling gives \(A=K=\beta=0\).
With \(Q_-=q(0)=8T\) and feasible spatial data the gap inequality
is exactly \(\Delta<T\), and is sharp by equality at \(T\).

For \(h=z^2+\alpha z^3+2\tau+6\alpha\tau z\) the deflated factor is
\(1+\alpha(a+b+x)\). It has \(g_z=-\alpha^2/U^2\). On a
positive-factor box the average in \eqref{eq:gap-law} is
\[
B=-\frac{\alpha^2}{(1+3\alpha c)^2-\alpha^2d^2}.
\]
At \((0,0)\) the polynomial has a nondegenerate double with
\(B_*=-\alpha^2\). The local splitting from Section~\ref{sec:geometry}
gives real branches before the event. A strictly negative \(B_*\)
for \(\alpha\ne0\) is compatible with collision.

For \(h=z^4+12\tau z^2+12\tau^2\) the origin is a quadruple
event with \(h_{zz}=0\). The repeated-node factor vanishes,
so the strict floor excludes this class. All three examples solve
the forward heat equation by direct differentiation. The local
events can be shifted in space and time. Their role
is to check the conditional real-pair argument.

\subsection{Comparison with contour acquisition}

A fixed separating contour also avoids internal-gap losses
\cite{RTTransfer}. To see the difference in required data, let
\(\Gamma\) be a simple positively oriented contour of length \(\ell\),
with the real segment at distance at least \(r>0\). Suppose it
encloses exactly the actual pair counted with multiplicity, with
the model \(F\) holomorphic across the closed interior and
\[
 |F|\ge m_\Gamma>\delta_0,\quad |h-F|\le\delta_0,\quad
 |F_z|\le M_\Gamma,\quad |h_z-F_z|\le\delta_1
 \qquad\text{on }\Gamma.
\]
Rouch\'e's theorem supplies the model count. The contour formula for
the true deflated logarithmic derivative is
\[
g(z)=\frac{1}{2\pi i}\int_\Gamma
        \frac{h_z(\zeta)/h(\zeta)}{\zeta-z}\,\dd\zeta .
\]
Indeed, the residues of the enclosed zeros subtract the logarithmic
derivative of their divisor from \(h_z/h\). The model formula
deflates by its own divisor, in contrast to the shared-node affine
subtraction in Section~\ref{sec:deflation}. The two logarithmic
derivatives on the contour differ by at most
\[
 Q_\Gamma=\frac{\delta_1+M_\Gamma\delta_0/m_\Gamma}
                   {m_\Gamma-\delta_0}.
\]
Differentiating the contour formulas gives
\(|g_z-(g_F)_z|\le \ell Q_\Gamma/(2\pi r^2)\) on the real segment.
Alternatively, the logarithm \(\log(h/F)\) is single-valued on a
collar of the contour and bounded there on \(\Gamma\) by
\(\nu=-\log(1-\delta_0/m_\Gamma)\). Integration by parts around
\(\Gamma\) moves the derivative to the Cauchy kernel. A subsequent
evaluation derivative gives
\(|g_z-(g_F)_z|\le \ell\nu/(\pi r^3)\).
The contour approach trades the uniform real-box factor floor for
complex boundary and buffer data and a paid cluster count.
Small root slopes alone settle neither method's relative cost.
